\documentclass[11pt,a4paper]{article}

\usepackage[utf8]{inputenc}
\usepackage[T1]{fontenc}
\usepackage{amsmath, amssymb, amsthm}
\usepackage{mathtools}
\usepackage{listings}
\usepackage{xcolor}
\usepackage{hyperref}
\usepackage[margin=1in]{geometry}
\usepackage{microtype}

\theoremstyle{plain}
\newtheorem{theorem}{Theorem}[section]
\newtheorem{lemma}[theorem]{Lemma}
\newtheorem{proposition}[theorem]{Proposition}
\newtheorem{corollary}[theorem]{Corollary}

\theoremstyle{definition}
\newtheorem{definition}[theorem]{Definition}
\newtheorem{example}[theorem]{Example}

\theoremstyle{remark}
\newtheorem{remark}[theorem]{Remark}

\lstdefinelanguage{Lean4}{
  morekeywords={theorem,lemma,def,structure,inductive,namespace,end,
    open,import,instance,where,fun,let,match,with,by,intro,apply,
    exact,sorry,have,show,refine,cases,rfl,Type,Prop,Sort},
  sensitive=true,
  morecomment=[l]{--},
  morecomment=[s]{/-}{-/},
  morestring=[b]"
}

\lstset{
  language=Lean4,
  basicstyle=\ttfamily\small,
  keywordstyle=\color{blue!70!black}\bfseries,
  commentstyle=\color{gray}\itshape,
  stringstyle=\color{red!70!black},
  showstringspaces=false,
  breaklines=true,
  frame=single,
  framesep=4pt,
  framerule=0.4pt,
  rulecolor=\color{gray!50},
  xleftmargin=0pt,
  xrightmargin=0pt
}

\newcommand{\Lan}{\mathrm{Lan}}
\newcommand{\Ran}{\mathrm{Ran}}
\newcommand{\Hom}{\mathrm{Hom}}
\newcommand{\id}{\mathrm{id}}
\newcommand{\Sm}{S_{m}}
\newcommand{\Ztwo}{\mathbb{Z}_{2}}
\newcommand{\fpath}[1]{\texttt{#1}}

\title{Arrow-Debreu, Arrow's Impossibility, and Schelling-Ising\\
       as Regimes of a Single Categorical Construction:\\
       A Lean 4 Formalization}

\author{Onyeka Obi\\
        MavenRain\\
        \texttt{claude@isurvivable.cv}}

\date{\today}

\begin{document}

\maketitle

\begin{abstract}
We give a Lean 4 formalization of a unification of three classical
aggregation theorems (Arrow-Debreu equilibrium, Arrow's Impossibility
Theorem, and the Schelling-Ising segregation model) as regimes of a
single categorical object: the left Kan extension of a choice rule along
the orbit projection of a symmetry group acting on a configuration
category.  Phase~1a wires Arrow's theorem into the framework via the
symmetric group action on preference profiles, with a strong
no-equivariant-constrained-SWF theorem proved end to end (the
\textit{anonymity} that the categorical $\Sm$ action encodes is shown to
be incompatible with Pareto plus Independence of Irrelevant Alternatives
for $m \geq 2$ voters and three or more alternatives).  Phase~1b realizes
Schelling-Ising via the $\Ztwo$ action on spin configurations, with the
symbolic bifurcation theorem (distinct ground states with equal
Hamiltonian witnessing $\Ztwo$-broken attractor pairing) proved.  The
formalization is built on top of \texttt{comp-cat-theory} (categorical
primitives), \texttt{kan-tactics} (term-mode proof primitives) and
\texttt{arrow-cat} (the Geanakoplos pivotal-voter proof of Arrow's
theorem), with no Mathlib dependency.  Two of the three regimes of the
trichotomy now have framework-load-bearing theorems proved end to end;
the third (Arrow-Debreu via Kakutani fixed point) and the analytic form
of the Schelling-Ising bifurcation remain as future work.
\end{abstract}

\section{Introduction}
\label{sec:introduction}

Three of the most cited results in mathematical economics and social
science live in disjoint mathematical communities.  The Arrow-Debreu
existence theorem for general equilibrium~\cite{arrow-debreu-1954} sits
in mathematical economics; Arrow's Impossibility
Theorem~\cite{arrow-1951} sits in social choice; and the Schelling
segregation model~\cite{schelling-1971}, formally a phase transition in
an Ising-type system, sits in statistical mechanics and the sociophysics
tradition that grew from it~\cite{galam-2012, bouchaud-2013}.  Each
result is a verdict on the same underlying problem: \emph{when does a
collection of local-level data (preferences, prices, choices, spins)
determine a coherent global-level outcome?}  Arrow-Debreu's verdict
is yes, under convexity.  Arrow's Impossibility's verdict is no, under
democratic aggregation axioms.  Schelling-Ising's verdict is multi-stable
yes, under spontaneous symmetry breaking.

The categorical observation we formalize is that all three verdicts are
properties of the same object.  Given a small category $C$ of micro-
configurations, a group $G$ acting on $C$ by symmetries, an outcome
category $D$, and a choice rule $F : C \to D$, the \emph{aggregation
map} is the left Kan extension $\Lan_p F$ along the orbit projection
$p : C \to C \sslash G$ into the action groupoid.  The three theorems
become three regimes of when and how this Kan extension is realized:
\begin{itemize}
  \item \textbf{Arrow-Debreu} (convex regime): $\Lan_p F$ exists and is
        uniquely realized, with the universal arrow giving a single
        equilibrium per orbit.
  \item \textbf{Arrow-Impossibility} (obstruction regime): no $\Lan_p F$
        exists satisfying $G$-equivariance ($=$ anonymity), unanimity
        ($=$ Pareto), and IIA-as-naturality simultaneously.
  \item \textbf{Schelling-Ising} (multi-stable regime): $\Lan_p F$ is
        realized by multiple non-isomorphic selectors, which bifurcate
        with a control parameter (inverse temperature) at a critical
        value.
\end{itemize}

The Brock-Durlauf statistical-mechanics-of-discrete-choice
program~\cite{brock-durlauf-2001} provides a $\beta$-parametrized family
$F_\beta$ of choice rules whose Kan extensions traverse all three
regimes as $\beta$ varies, suggesting that the trichotomy is a phase
diagram rather than three independent classification results.  Topological
social choice theory~\cite{chichilnisky-1980} reformulates Arrow's
Impossibility as the non-contractibility of an appropriate equivariant
configuration space; the obstruction regime above is the categorical
shadow of that topological obstruction.

\paragraph{Contribution.}
We present a Lean~4 formalization of this unification.  The contribution
is fourfold.

First, we lay out the categorical scaffolding (\fpath{ConfigSpace},
\fpath{ChoiceRule}, \fpath{Aggregation}, \fpath{Regimes},
\fpath{Trichotomy}) and the auxiliary infrastructure
(\fpath{Discrete} categories, \fpath{SymmetricGroup} and \fpath{Z2Group}
as instances of \fpath{SymmetryGroup}, two flavors of
\fpath{Functor} extensionality including a load-bearing
\fpath{Functor.ext\_pointwise} lemma that handles obj-equality-up-to-
pointwise-rewrite cases).  All of this is built on
\texttt{comp-cat-theory}~\cite{compcat-theory}, with no Mathlib
dependency.

Second, we wire Arrow's Impossibility Theorem into the framework via the
$\Sm$ action on \fpath{Profile m \(\alpha\)} (Section~\ref{sec:phase1a}).
The headline theorem
\fpath{no\_equivariant\_constrained\_swf} states that no social welfare
function over $m \geq 2$ voters and three or more alternatives can
simultaneously be $\Sm$-equivariant and satisfy Pareto plus IIA.  Its
proof chain (\fpath{equivariant\_dictator\_transfer},
\fpath{all\_voters\_dictator\_under\_equivariance},
\fpath{no\_two\_dictators}) makes the framework's
\fpath{GAction} structure non-decorative: the anonymity that the
categorical $\Sm$-action encodes is the very obstruction.  arrow-cat's
classical statement~\cite{arrow-cat} permits a dictator solution; the
framework's anonymity requirement rules that out.

Third, we realize the Schelling-Ising regime via the $\Ztwo$ action on
spin configurations (Section~\ref{sec:phase1b}).  At the categorical
level, \fpath{spinConfigAction : GAction Z2Group (SpinConfigCat n)} is
fully proved using \fpath{Functor.ext\_pointwise}.  At the
phase-transition level, we define magnetization and the mean-field
ferromagnetic Hamiltonian (integer-valued, $J = 1$), prove their $\Ztwo$
inversion and invariance laws, and then prove a \emph{symbolic} form of
the bifurcation theorem: \fpath{schelling\_ising\_z2\_degeneracy}
witnesses that for $n \geq 1$ the Hamiltonian admits at least two
distinct configurations of equal energy.  This is the structural content
of bifurcation (distinct attractors of the $\Ztwo$-equivariant
dynamics) captured within an Int-level foundation that needs no
real-number machinery.

Fourth, we identify what remains.  The Arrow-Debreu regime
(Phase~2) awaits a categorical Kakutani fixed-point argument.  The
analytic Schelling-Ising bifurcation theorem (one component for
$\beta < \beta_c$, two for $\beta > \beta_c$, with $\beta_c$ a specific
real number) awaits a real-valued analysis layer, most naturally
Mathlib.  Four cast-tower \fpath{HEq} sorries in
\fpath{Aggregation.lean}'s orbit-groupoid construction remain as
technical debt; they are orthogonal to the headline theorems but
prevent the orbit-groupoid laws from being fully checked.

\paragraph{Outline.}
Section~\ref{sec:object} states the central categorical object and the
three regimes.  Section~\ref{sec:foundations} sketches the foundational
infrastructure (discrete categories, group actions, the two Functor
extensionality lemmas).  Section~\ref{sec:phase1a} presents the
Arrow-Impossibility wiring.  Section~\ref{sec:phase1b} presents the
Schelling-Ising wiring and the symbolic bifurcation theorem.
Section~\ref{sec:discussion} discusses related work, the framework's
load-bearing properties, and the path to closing out the remaining
regimes and sorries.

\section{The Central Object and the Three Regimes}
\label{sec:object}

\begin{definition}[Configuration space with symmetry]
A \emph{configuration space with symmetry} is a triple $(C, G, \alpha)$
where $C$ is a small category, $G$ is a group, and $\alpha : G \to
\mathrm{End}(C)$ is a (covariant) action of $G$ on $C$ by endofunctors,
that is, $\alpha(e) = \id_C$ and $\alpha(gh) = \alpha(h) \circ
\alpha(g)$.\footnote{We use the diagrammatic ordering convention
$\alpha(gh) = \alpha(h) \circ \alpha(g)$ throughout, matching the
\texttt{comp-cat-theory} convention $F \circ G = F \,\fpath{⋙}\, G$
for ``apply $F$ first, then $G$.''}
\end{definition}

\begin{definition}[Choice rule]
A \emph{choice rule} on $(C, G, \alpha)$ valued in an outcome category
$D$ is a functor $F : C \to D$.  A \emph{$\beta$-family of choice
rules} is a function $\beta \mapsto F_\beta$ from a parameter type
$\beta$ to choice rules $C \to D$.
\end{definition}

\begin{definition}[Aggregation]
The \emph{aggregation} of a choice rule $F$ along the symmetry
$\alpha$ is the left Kan extension
\[
  \mathrm{Aggregation}(\alpha, F) = \Lan_{p_\alpha} F
\]
of $F$ along the orbit projection $p_\alpha : C \to C \sslash G$ into
the action groupoid.  $C \sslash G$ has the same objects as $C$ and
morphisms $X \to Y$ given by pairs $(g, f)$ with $g \in G$ and
$f : X \to \alpha(g)(Y)$ a morphism in $C$.
\end{definition}

The three regimes are then properties of when and how
$\Lan_{p_\alpha} F$ is realized.

\begin{description}
  \item[Arrow-Debreu regime.] $\Lan_{p_\alpha} F$ is realized and the
        universal cocone has a unique vertex per orbit.  Sufficient
        condition: $F$ is valued in a convex subcategory of $D$ and the
        action $\alpha$ is by group-of-automorphism-preserving functors.
        The classical Arrow-Debreu existence theorem~\cite{arrow-debreu-1954}
        is the special case where $C$ is a discrete category of
        consumption-bundle endowments and $D$ is a category of price
        vectors with morphisms given by the price-adjustment dynamics.
  \item[Arrow-Impossibility regime.] No $\Lan_{p_\alpha} F$ exists that
        simultaneously satisfies $G$-equivariance, unanimity, and
        IIA-as-naturality.  The obstruction is a class in an equivariant
        cohomology group (Chichilnisky's topological proof of Arrow's
        theorem~\cite{chichilnisky-1980} is the topological version of
        this statement).  Section~\ref{sec:phase1a} formalizes this for
        $C =$ discrete category of preference profiles, $G = \Sm$,
        $D =$ discrete category of strict preferences.
  \item[Schelling-Ising regime.] $\Lan_{p_\alpha} F$ is realized but
        multi-valued in the stack-theoretic sense: the universal cocone
        is realized by multiple non-isomorphic selectors.  Phase
        transitions are bifurcations in the homotopy type of the
        selector space as a parameter (e.g.\ Brock-Durlauf's inverse
        temperature $\beta$) varies.  Section~\ref{sec:phase1b}
        formalizes the symbolic form of this for $C =$ discrete
        category of spin configurations, $G = \Ztwo$, $D =$ outcome
        space; the analytic form remains future work.
\end{description}

\begin{remark}[Symmetry as the unifying thread]
A common thread across the three regimes is the question of which
symmetries of the local-level data survive aggregation.  Arrow-Debreu
preserves Pareto and budget-balance symmetries.  Arrow's Impossibility
says certain anonymity-style symmetries cannot be jointly preserved by
any aggregation rule.  Spontaneous symmetry breaking in Schelling-Ising
is the statement that the dynamical realization breaks a $\Ztwo$
symmetry that the energy functional preserves.  Reformulating these
as statements about whether $\Lan_{p_\alpha} F$ admits a $G$-equivariant
realization makes the symmetry content of each theorem explicit.
\end{remark}

\section{Foundations}
\label{sec:foundations}

\textit{Section to be expanded: discrete categories, group actions on
categories, the orbit groupoid construction with explicit cast machinery
through act\_one and act\_mul, and the two Functor extensionality
lemmas.  The latter, Functor.ext (HEq form) and
Functor.ext\_pointwise (pointwise form), are load-bearing for the
Phase~1b act\_mul law and would be of independent interest as additions
to comp-cat-theory.}

\section{Phase 1a: Arrow-Impossibility via $\Sm$ equivariance}
\label{sec:phase1a}

We instantiate the framework for Arrow's setting: $C$ is the discrete
category on \fpath{Profile m \(\alpha\)} (functions $\mathrm{Fin}\, m
\to \mathrm{StrictPref}\, \alpha$), $G = \Sm$ (the symmetric group on
$m$ voter indices), $D$ is the discrete category on
$\mathrm{StrictPref}\, \alpha$.  The $\Sm$ action on profiles is the
pointwise reindexing
$(\sigma \cdot p)(i) = p(\sigma(i))$.

\subsection{The categorical embedding}

\begin{lstlisting}
abbrev ProfileCat (m : Nat) (α : Type u) : Type u :=
  Discrete (Profile m α)

def actByPerm {m : Nat} {α : Type u} (σ : Perm m) :
    ProfileCat m α ⥤ ProfileCat m α :=
  discreteFunctor (fun p => fun i => p (σ.toFun i))

def profileAction (m : Nat) (α : Type u) :
    GAction (SymmetricGroup m) (ProfileCat m α) where
  act := actByPerm
  act_one := /- via Functor.ext + funext + heq_of_eq + match -/
  act_mul := /- same recipe -/
\end{lstlisting}

The \fpath{act\_one} and \fpath{act\_mul} laws close via the
\fpath{Functor.ext} recipe, exploiting structure-eta on the \fpath{Discrete}
wrapper plus definitional proof irrelevance on the
\fpath{map\_id} and \fpath{map\_comp} fields.

\subsection{The headline theorem}

The framework treats a social welfare function as a choice rule
\fpath{SWFasChoiceRule f : ProfileCat m \(\alpha\) ⥤ Discrete
(StrictPref \(\alpha\))} (a \fpath{discreteFunctor} applied to the
underlying \fpath{f : SWF m \(\alpha\)}).  The headline theorem is:

\begin{theorem}[no\_equivariant\_constrained\_swf]
\label{thm:no-equivariant-constrained-swf}
For $m \geq 2$ voters, a nonempty profile space, and three or more
alternatives, no SWF $f : \fpath{SWF m \(\alpha\)}$ is simultaneously
$\Sm$-equivariant and satisfies Pareto plus IIA.
\end{theorem}

\begin{proof}[Proof sketch]
The proof chain uses three intermediate lemmas, all proved in the
formalization:
\begin{enumerate}
  \item \fpath{equivariant\_dictator\_transfer}: under an
        $\Sm$-equivariant SWF, the dictator role transfers between any
        two voters connected by a permutation.  This is the
        load-bearing lemma: without equivariance the dictator from
        Arrow's theorem~\cite{arrow-cat} stays fixed; with equivariance
        it propagates.
  \item \fpath{all\_voters\_dictator\_under\_equivariance}: combining
        Arrow's theorem with the transfer lemma plus transitivity of
        $\Sm$ on $\mathrm{Fin}\, m$, every voter is a dictator under
        equivariance.  Transitivity is witnessed inline by
        \texttt{ArrowCat.swapMap} (a swap permutation lifted to
        $\fpath{Perm m}$ as an involutive bijection).
  \item \fpath{no\_two\_dictators}: two distinct voters cannot both be
        dictators.  Proof builds a disagreement profile using
        \texttt{ArrowCat.StrictPref.moveBToBottom}: voter $k_1$ has $b$
        moved to the bottom (so $a > b$), voter $k_2$ has $a$ moved to
        the bottom (so $b > a$).  Both dictators then force
        $(f p).\mathrm{pref}\, a\, b$ and $(f p).\mathrm{pref}\, b\, a$,
        contradicting \fpath{StrictPref.asym}.
\end{enumerate}
For $m \geq 2$ the index $0, 1 \in \mathrm{Fin}\, m$ are distinct
voters, and the combination of (2) and (3) closes the impossibility.
\end{proof}

\begin{remark}[Framework load-bearing]
Theorem~\ref{thm:no-equivariant-constrained-swf} is stronger than the
classical Arrow statement: arrow-cat's \fpath{arrow} permits a dictator
solution to Pareto plus IIA, but the framework's
$\Sm$-equivariance requirement turns the dictator into a contradiction
for $m \geq 2$.  The categorical $\Sm$ action does work that the
classical statement does not; the framework is non-decorative.
\end{remark}

\section{Phase 1b: Schelling-Ising via $\Ztwo$ symmetry}
\label{sec:phase1b}

\textit{Section to be expanded.  We instantiate the framework for
Schelling-Ising at the symbolic level: $C$ is the discrete category
on $\fpath{SpinConfig n}$ (functions $\fpath{Fin n} \to \fpath{Spin}$),
$G = \Ztwo$, the action is pointwise spin flip.  The headline pieces
are:
\begin{itemize}
  \item \fpath{spinConfigAction : GAction Z2Group (SpinConfigCat n)},
        proved using \fpath{Functor.ext\_pointwise} on the
        \fpath{act\_mul (.g, .g)} case where the obj equality requires
        \fpath{SpinConfig.flip\_flip} rather than rfl.
  \item \fpath{Magnetization\_flip}: flipping all spins negates the
        magnetization.  The $\Ztwo$ inversion law for the order
        parameter.
  \item \fpath{Hamiltonian\_flip}: the mean-field ferromagnetic Ising
        Hamiltonian $H(c) = -(M(c))^2$ is $\Ztwo$-invariant.  Forces
        any Hamiltonian-based choice rule to inherit the symmetry.
  \item \fpath{schelling\_ising\_z2\_degeneracy}: for $n \geq 1$, the
        Hamiltonian admits at least two distinct configurations of
        equal energy, namely \fpath{upConfig n} and \fpath{downConfig
        n}.  The \emph{symbolic} form of bifurcation, capturing the
        $\Ztwo$-broken attractor pairing without invoking real-valued
        machinery.
\end{itemize}
The analytic form of the bifurcation theorem (one component for
$\beta < \beta_c$, two for $\beta > \beta_c$, with $\beta_c$ a specific
real number arising from $m = \tanh(\beta J m)$) requires real-valued
infrastructure beyond Int and remains future work; the most natural
path adds Mathlib as a dependency.}

\section{Discussion}
\label{sec:discussion}

\textit{Section to be expanded.  Topics: relation to topological social
choice~\cite{chichilnisky-1980} and equivariant cohomology of the
configuration space; relation to Brock-Durlauf's statistical-mechanics-
of-discrete-choice~\cite{brock-durlauf-2001} as the $\beta$-family that
traverses the regimes; the open Arrow-Debreu regime (Phase 2) and a
sketch of a categorical Kakutani argument; the open analytic Schelling-
Ising bifurcation theorem and the path through Mathlib; the open
orbit-groupoid \fpath{HEq} sorries in \fpath{Aggregation.lean} and the
prospect of a unified cast-manipulation tactic.}

\section*{Reproducibility and Code Availability}

The formalization is available at
\url{https://github.com/MavenRain/unified-aggregation-theory} under a
dual MIT/Apache-2.0 license.  The release tags
\fpath{v0.1.0-phase-1a} and \fpath{v0.1.2-phase-1b-bifurcation-symbolic}
correspond to the Phase~1a and Phase~1b symbolic completion states
described above.  API documentation is auto-generated via doc-gen4 at
\url{https://MavenRain.github.io/unified-aggregation-theory/}.

Dependencies are \texttt{comp-cat-theory}, \texttt{kan-tactics}, and
\texttt{arrow-cat}, all sibling-directory \texttt{lake require}'d.  The
Lean toolchain is \texttt{leanprover/lean4:v4.30.0-rc2}; no Mathlib
dependency.

\begin{thebibliography}{99}

\bibitem{arrow-1951}
K.\ J.\ Arrow.
\textit{Social Choice and Individual Values}.
Wiley, 1951.

\bibitem{arrow-debreu-1954}
K.\ J.\ Arrow and G.\ Debreu.
Existence of an equilibrium for a competitive economy.
\textit{Econometrica}, 22(3):265--290, 1954.

\bibitem{schelling-1971}
T.\ C.\ Schelling.
Dynamic models of segregation.
\textit{Journal of Mathematical Sociology}, 1(2):143--186, 1971.

\bibitem{chichilnisky-1980}
G.\ Chichilnisky.
Social choice and the topology of spaces of preferences.
\textit{Advances in Mathematics}, 37(2):165--176, 1980.

\bibitem{brock-durlauf-2001}
W.\ A.\ Brock and S.\ N.\ Durlauf.
Discrete choice with social interactions.
\textit{Review of Economic Studies}, 68(2):235--260, 2001.

\bibitem{galam-2012}
S.\ Galam.
\textit{Sociophysics: A Physicist's Modeling of Psycho-political
Phenomena}.
Springer, 2012.

\bibitem{bouchaud-2013}
J.-P.\ Bouchaud.
Crises and collective socio-economic phenomena: simple models and
challenges.
\textit{Journal of Statistical Physics}, 151(3-4):567--606, 2013.

\bibitem{compcat-theory}
O.\ Obi.
comp-cat-theory: A Lean 4 reservoir library for compositional category
theory.
\url{https://github.com/MavenRain/comp-cat-theory}.

\bibitem{arrow-cat}
O.\ Obi.
arrow-cat: A Lean 4 formalization of Arrow's Impossibility Theorem via
the Geanakoplos pivotal-voter argument.
\url{https://github.com/MavenRain/arrow-cat}.

\end{thebibliography}

\end{document}
